\chapter{Images in Python}

Every photo filter on your phone, every face detector, and every
self-driving car's camera does the same basic thing: it treats a picture as
a big grid of numbers and does math on them. In this chapter, you will
make images from scratch with Python and take real photos apart.

As you saw in the camera chapter, a color image is a grid of pixels, each
with a red, green, and blue value from 0 to 255. In Python, we store that
as a three-dimensional \index{NumPy array}NumPy array of 8-bit integers
(type \texttt{uint8}, which holds exactly the values 0 to 255).

If you have an \index{RGB}RGB image that is 480 \index{pixel}pixels tall and 640 pixels wide,
you will need a $480 \times 640 \times 3$ NumPy array.

A separate library, \index{imageio}imageio, does the file handling. (Install it with
\texttt{pip3 install imageio} if you need to.) It
\begin{itemize}
\item Reads an image file (like JPEG files) and creates a NumPy array.
\item Writes a NumPy array to a file in standard image formats
\end{itemize}

Let us create a simple Python program that creates a file containing an
all-black image that is 640 pixels wide and 480 pixels tall. Create a
file called \filename{create\_image.py}:

\begin{minted}{python}
import numpy as np
import imageio.v3 as iio
import sys

# Check command-line arguments
if len(sys.argv) < 2:
    print(f"Usage {sys.argv[0]} <outfile>")
    sys.exit(1)

# Constants
IMAGE_WIDTH = 640
IMAGE_HEIGHT = 480

# Create an array of zeros
image = np.zeros((IMAGE_HEIGHT, IMAGE_WIDTH, 3), dtype=np.uint8)

# Write the array to the file
iio.imwrite(sys.argv[1], image)
\end{minted}

To run this, you supply the name of the file you want to create as a
\index{command-line argument}\newterm{command-line argument}; Python puts it
in \pyvar{sys.argv[1]}. The extension (like .png or .jpeg) will tell imageio
what format you want written. Run it now:

\begin{minted}{python}
python3 create_image.py blackness.png
\end{minted}

Open the image to confirm that it is 640 pixels wide, 480 pixels tall, and completely black.

\section{Adding color}\index{color (RGB)}

\textbf{Worked example.} What color is a pixel with the value
$(255, 255, 0)$?

\textbf{Solution.} Full red, full green, and no blue. As you learned in the
eye chapter, red and green light together look yellow.


Now, let us walk through the image, pixel-by-pixel, adding some
red. We will gradually increase the red from 0 on the left to 255 on
the right.

\begin{minted}{python}
import numpy as np
import imageio.v3 as iio
import sys

# Check command-line arguments
if len(sys.argv) < 2:
    print(f"Usage {sys.argv[0]} <outfile>")
    sys.exit(1)

# Constants
IMAGE_WIDTH = 640
IMAGE_HEIGHT = 480

# Create an array of zeros
image = np.zeros((IMAGE_HEIGHT, IMAGE_WIDTH, 3), dtype=np.uint8)

for col in range(IMAGE_WIDTH):

    # Red goes from 0 to 255 (left to right)
    r = int(col * 255.0 / IMAGE_WIDTH)

    # Update all the pixels in that column
    for row in range(IMAGE_HEIGHT):
        # Set the red pixel
        image[row, col, 0] = r

# Write the array to the file
iio.imwrite(sys.argv[1], image)
\end{minted}

When you run the function to create a new image, it will be a fade
from black to red as you move from left to right:

\includegraphics[width=0.4\linewidth]{red.png}

Now, inside the inner loop, update the blue channel so that it goes
from zero at the top to 255 at the bottom:

\begin{minted}{python}
    # Update all the pixels in that column
    for row in range(IMAGE_HEIGHT):

        # Update the red channel
        image[row,col,0] = r
        
        # Blue goes from 0 to 255 (top to bottom)
        b = int(row * 255.0 / IMAGE_HEIGHT)
        image[row,col,2] = b

iio.imwrite(sys.argv[1], image)
\end{minted}

When you run the program again, you will see the color fades from
black to blue as you go down the left side. As you go down the right
side, the color fades from red to magenta.

\includegraphics[width=0.4\linewidth]{rb.png}

Notice that red and blue with no green looks magenta to your eye.

Next, let us add some stripes of green:

\begin{minted}{python}
import numpy as np
import imageio.v3 as iio
import sys

# Check command line arguments
if len(sys.argv) < 2:
    print(f"Usage {sys.argv[0]} <outfile>")
    sys.exit(1)

# Constants
IMAGE_WIDTH = 640
IMAGE_HEIGHT = 480
STRIPE_WIDTH = 40
pattern_width = STRIPE_WIDTH * 2

# Create an image of all zeros
image = np.zeros((IMAGE_HEIGHT, IMAGE_WIDTH, 3), dtype=np.uint8)

# Step from left to right
for col in range(IMAGE_WIDTH):

    # Red goes from 0 to 255 (left to right)
    r = int(col * 255.0 / IMAGE_WIDTH)

    # Should I add green to this column?
    should_green = col % pattern_width > STRIPE_WIDTH

    # Update all the pixels in that column
    for row in range(IMAGE_HEIGHT):

        # Update the red channel
        image[row,col,0] = r

        # Should I add green to this pixel?
        if should_green:
            image[row,col,1] = 255

        # Blue goes from 0 to 255 (top to bottom)
        b = int(row * 255.0 / IMAGE_HEIGHT)
        image[row,col,2] = b

iio.imwrite(sys.argv[1], image)
\end{minted}

When you run this version, you will see the previous image in every
other stripe.  In the other half, you will see that green fades to cyan
down the left side, and yellow fades to white down the right side.

\includegraphics[width=0.4\linewidth]{rgb.png}

\section{Using an existing image}

\index{image file formats}imageio can also be used to read in any common image file
format. Let us read in an image and save each of the red, green, and blue
channels out as its own image.

Create a new file called \filename{separate\_image.py}:

\begin{minted}{python}
import numpy as np
import imageio.v3 as iio
import sys
import os

# Check command line arguments
if len(sys.argv) < 2:
    print(f"Usage {sys.argv[0]} <infile>")
    sys.exit(1)

# Read the image
path = sys.argv[1]
image = iio.imread(path)

# What is the filename?
filename = os.path.basename(path)

# What is the shape of the array?
original_shape = image.shape

# Log it
print(f"Shape of {filename}:{original_shape}")

# Names of the colors for the filenames
colors = ['red','green','blue']

# Step through each of the colors
for i in range(3):

    # Create a new image
    newimage = np.zeros(original_shape, dtype=np.uint8)

    # Copy one channel
    newimage[:,:,i] = image[:,:,i]

    # Save to a file
    new_filename = f"{colors[i]}_{filename}"
    print(f"Writing {new_filename}")
    iio.imwrite(new_filename, newimage)
\end{minted}

Now, you can run the program with any common RGB image type:

\begin{minted}{python}
python3 separate_image.py dog.jpg
\end{minted}

This will create three images: \filename{red\_dog.jpg},
\filename{green\_dog.jpg}, and \filename{blue\_dog.jpg}. Look at them side
by side. A red sweater is bright in the red image and dark in the others;
the sky is brightest in the blue one.

\section{Faster with NumPy}
The loops above visit all 307,200 pixels one at a time, which is slow.
NumPy can do the same work on a whole array at once. This one line makes
the left-to-right red fade without any loops:
\begin{minted}{python}
image[:, :, 0] = np.linspace(0, 255, IMAGE_WIDTH).astype(np.uint8)
\end{minted}
NumPy copies the row of 640 values down every one of the 480 rows. Real
image-processing code is written this way, and it can be hundreds of times
faster.

\begin{Exercise}[title={Grayscale}, label=img_gray]
Write a program that reads a color image and writes a grayscale version.
A common formula for how bright a pixel looks is
$0.299R + 0.587G + 0.114B$. Why is green weighted most heavily?
\end{Exercise}
\begin{Answer}[ref=img_gray]
\begin{minted}{python}
import numpy as np
import imageio.v3 as iio
import sys

image = iio.imread(sys.argv[1]).astype(float)
gray = 0.299 * image[:, :, 0] + 0.587 * image[:, :, 1] + 0.114 * image[:, :, 2]
iio.imwrite(sys.argv[2], gray.astype(np.uint8))
\end{minted}
Green gets the most weight because our eyes are most sensitive to green
light, so green contributes the most to how bright a color looks. (The
weights add up to 1, so white stays 255.)
\end{Answer}

\begin{Exercise}[title={White Pixel}, label=img1]
What RGB value makes a pixel white?
\end{Exercise}
\begin{Answer}[ref=img1]
$(255, 255, 255)$, with full red, full green, and full blue.
\end{Answer}
